% ======================================================================
%  LATEX FIGURE PRACTICE PACK  ---  10 graded exercises + full solutions
%  ------------------------------------------------------------------
%  HOW TO USE THIS FILE
%    1. Compile TWICE:  pdflatex figure_practice.tex
%    2. PRINT (or keep open) only PART I -- the target figures.
%    3. Open a BLANK .tex file, start a timer, reproduce each target
%       from scratch. Do NOT scroll to Part II.
%    4. Only after your attempt compiles, check PART II for the solution
%       and the marking notes.
%
%  All images are the placeholders shipped with the `mwe' package:
%    example-image        example-image-a    example-image-b
%    example-image-c      example-image-16x9 example-image-golden
%  In the real exam, swap these for the supplied filenames.
% ======================================================================
\documentclass[11pt]{article}

\usepackage[a4paper,margin=1in]{geometry}
\usepackage[T1]{fontenc}
\usepackage[utf8]{inputenc}

\usepackage{amsmath,amssymb}
\usepackage[dvipsnames]{xcolor}
\usepackage{graphicx}
\usepackage{mwe}              % supplies the example-image-* placeholders
\usepackage{subcaption}       % subfigures + \subref (loads caption)
\usepackage{wrapfig}          % text-wrapped figures
\usepackage{rotating}         % sidewaysfigure
\usepackage{float}            % [H] exact placement
\usepackage{array}
\usepackage{tabularx}
\usepackage{booktabs}
\usepackage{listings}
\usepackage{enumitem}
\usepackage{hyperref}
\hypersetup{colorlinks=true, linkcolor=RoyalBlue, urlcolor=magenta}

\lstdefinestyle{tex}{
  language=[LaTeX]TeX, basicstyle=\ttfamily\footnotesize,
  keywordstyle=\color{RoyalBlue}\bfseries,
  commentstyle=\color{gray}\itshape,
  backgroundcolor=\color{black!3}, frame=single, framerule=0.3pt,
  breaklines=true, columns=fullflexible, showstringspaces=false,
  keepspaces=true,
  morekeywords={includegraphics,subfigure,minipage,wrapfigure,sidewaysfigure,
  rotatebox,caption,captionof,subref,label,centering,hfill,linewidth,
  textwidth,resizebox,fbox}}

\newcommand{\exinfo}[3]{%
  \par\noindent
  \colorbox{RoyalBlue!12}{\parbox{0.97\textwidth}{\small
  \textbf{Difficulty:} #1 \qquad \textbf{Target time:} #2 \\
  \textbf{Tests:} #3}}\par\medskip}

% filler prose for the text-wrapping exercises
\newcommand{\filler}{The transform method replaces differentiation in the
time domain with multiplication in the frequency domain, which turns an
awkward differential relation into ordinary algebra. Once the algebraic
problem has been solved, the inverse transform carries the result back to
the time domain. This detour is worthwhile because algebraic manipulation
is mechanical, whereas solving differential equations directly requires
case-by-case ingenuity.}

\setlength{\parindent}{0pt}
\setlength{\parskip}{4pt}

\title{\textbf{\LaTeX{} Figure Practice Pack}\\[2pt]
  \large 10 Graded Exercises Modelled on the Lab-Exam Papers}
\author{Reproduce Part I from scratch --- check Part II afterwards}
\date{\today}

\begin{document}
\maketitle

\begin{center}
\fbox{\parbox{0.92\textwidth}{\small\centering
\textbf{Rules of engagement.} Print Part I only. Start from a blank file with
your own preamble. Time each attempt. A figure ``passes'' only if it
\emph{compiles} and matches the target structurally --- panels in the right
arrangement, captions and sub-captions present, labels resolving in the text.}}
\end{center}

\tableofcontents
\newpage

% ======================================================================
\part*{PART I --- Target Figures}
\addcontentsline{toc}{section}{PART I --- Target Figures}
\section*{Reproduce each of the following}
\addcontentsline{toc}{subsection}{Exercises 1--10}

% ----------------------------------------------------------------------
\subsection*{Exercise 1 --- Warm-up}
\exinfo{$\star$}{3 minutes}{\texttt{figure} float, \texttt{centering},
width scaling, \texttt{caption} + \texttt{label}, cross-reference}

Figure~\ref{fig:warm} shows the reference waveform used throughout.

\begin{figure}[H]
  \centering
  \includegraphics[width=0.45\textwidth]{example-image-a}
  \caption{The reference waveform.}
  \label{fig:warm}
\end{figure}

As shown in Figure~\ref{fig:warm}, the response settles quickly.

% ----------------------------------------------------------------------
\subsection*{Exercise 2 --- Two panels, no sub-captions}
\exinfo{$\star\star$}{5 minutes}{\texttt{minipage} side by side,
\texttt{hfill}, one shared caption}

\begin{figure}[H]
  \centering
  \begin{minipage}{0.45\textwidth}
    \centering
    \includegraphics[width=\linewidth]{example-image-a}
  \end{minipage}\hfill
  \begin{minipage}{0.45\textwidth}
    \centering
    \includegraphics[width=\linewidth]{example-image-b}
  \end{minipage}
  \caption{Input signal (left) and filtered output (right).}
  \label{fig:pair}
\end{figure}

% ----------------------------------------------------------------------
\subsection*{Exercise 3 --- Three sub-captioned panels}
\exinfo{$\star\star$}{6 minutes}{\texttt{subfigure}, individual
sub-captions and labels, \texttt{ref} vs \texttt{subref}}

\begin{figure}[H]
  \centering
  \begin{subfigure}[b]{0.3\textwidth}
    \centering
    \includegraphics[width=\linewidth]{example-image-a}
    \caption{Under-damped}
    \label{fig:under}
  \end{subfigure}\hfill
  \begin{subfigure}[b]{0.3\textwidth}
    \centering
    \includegraphics[width=\linewidth]{example-image-b}
    \caption{Critically damped}
    \label{fig:crit}
  \end{subfigure}\hfill
  \begin{subfigure}[b]{0.3\textwidth}
    \centering
    \includegraphics[width=\linewidth]{example-image-c}
    \caption{Over-damped}
    \label{fig:over}
  \end{subfigure}
  \caption{Three damping regimes of a second-order system.}
  \label{fig:damping}
\end{figure}

Panel~\subref{fig:crit} of Figure~\ref{fig:damping} settles fastest, while
Figure~\ref{fig:under} overshoots.

% ----------------------------------------------------------------------
\newpage
\subsection*{Exercise 4 --- Nine panels in labelled groups}
\exinfo{$\star\star\star\star$}{12 minutes}{the Solar-System grid ---
group headings inside a float, row breaks, nine sub-labels}

Figure~\ref{fig:grid} shows all nine bodies.

\begin{figure}[H]
  \centering

  {\large\textbf{Star}}\par\medskip
  \begin{subfigure}[b]{0.3\textwidth}
    \centering
    \includegraphics[width=0.6\linewidth]{example-image-a}
    \caption{The Sun}\label{fig:sun}
  \end{subfigure}

  \bigskip
  {\large\textbf{Inner Planets}}\par\medskip
  \begin{subfigure}[b]{0.22\textwidth}
    \centering\includegraphics[width=\linewidth]{example-image-b}
    \caption{Mercury}\label{fig:mercury}
  \end{subfigure}\hfill
  \begin{subfigure}[b]{0.22\textwidth}
    \centering\includegraphics[width=\linewidth]{example-image-c}
    \caption{Venus}\label{fig:venus}
  \end{subfigure}\hfill
  \begin{subfigure}[b]{0.22\textwidth}
    \centering\includegraphics[width=\linewidth]{example-image-a}
    \caption{Earth}\label{fig:earth}
  \end{subfigure}\hfill
  \begin{subfigure}[b]{0.22\textwidth}
    \centering\includegraphics[width=\linewidth]{example-image-b}
    \caption{Mars}\label{fig:mars}
  \end{subfigure}

  \bigskip
  {\large\textbf{Outer Planets}}\par\medskip
  \begin{subfigure}[b]{0.22\textwidth}
    \centering\includegraphics[width=\linewidth]{example-image-c}
    \caption{Jupiter}\label{fig:jupiter}
  \end{subfigure}\hfill
  \begin{subfigure}[b]{0.22\textwidth}
    \centering\includegraphics[width=\linewidth]{example-image-a}
    \caption{Saturn}\label{fig:saturn}
  \end{subfigure}\hfill
  \begin{subfigure}[b]{0.22\textwidth}
    \centering\includegraphics[width=\linewidth]{example-image-b}
    \caption{Uranus}\label{fig:uranus}
  \end{subfigure}\hfill
  \begin{subfigure}[b]{0.22\textwidth}
    \centering\includegraphics[width=\linewidth]{example-image-c}
    \caption{Neptune}\label{fig:neptune}
  \end{subfigure}

  \caption{The star and the eight planets, arranged into the inner and
  outer groups. The displayed sizes and distances are not to scale.}
  \label{fig:grid}
\end{figure}

Compare Mars in Figure~\ref{fig:mars} with Jupiter in Figure~\ref{fig:jupiter}.

% ----------------------------------------------------------------------
\newpage
\subsection*{Exercise 5 --- Text wrapped around a figure}
\exinfo{$\star\star\star$}{7 minutes}{\texttt{wrapfigure}, side selection,
reserved width, caption inside a wrapped float}

\begin{wrapfigure}{r}{0.38\textwidth}
  \centering
  \includegraphics[width=\linewidth]{example-image-a}
  \caption{A wrapped figure on the right.}
  \label{fig:wrap}
\end{wrapfigure}

Figure~\ref{fig:wrap} sits on the right while this paragraph flows around
it. \filler

The text continues past the bottom of the float and returns to the full
measure once the figure has been cleared. \filler

% ----------------------------------------------------------------------
\subsection*{Exercise 6 --- Rotation in three flavours}
\exinfo{$\star\star\star$}{7 minutes}{\texttt{angle=} inside
\texttt{includegraphics}, \texttt{rotatebox} on arbitrary content}

\begin{figure}[H]
  \centering
  \begin{subfigure}[b]{0.28\textwidth}
    \centering
    \includegraphics[width=0.7\linewidth]{example-image-a}
    \caption{Upright}
    \label{fig:rot0}
  \end{subfigure}\hfill
  \begin{subfigure}[b]{0.28\textwidth}
    \centering
    \includegraphics[width=0.7\linewidth, angle=90]{example-image-a}
    \caption{Rotated $90^{\circ}$}
    \label{fig:rot90}
  \end{subfigure}\hfill
  \begin{subfigure}[b]{0.28\textwidth}
    \centering
    \rotatebox{30}{\fbox{Rotated text}}
    \caption{\texttt{rotatebox} on text}
    \label{fig:rottext}
  \end{subfigure}
  \caption{Three ways to rotate content.}
  \label{fig:rotations}
\end{figure}

% ----------------------------------------------------------------------
\newpage
\subsection*{Exercise 7 --- A $2 \times 3$ grid}
\exinfo{$\star\star\star$}{9 minutes}{two rows of three panels, blank line
as a row break, consistent widths}

\begin{figure}[H]
  \centering
  \begin{subfigure}[b]{0.3\textwidth}
    \centering\includegraphics[width=\linewidth]{example-image-a}
    \caption{Stage 1}\label{fig:s1}
  \end{subfigure}\hfill
  \begin{subfigure}[b]{0.3\textwidth}
    \centering\includegraphics[width=\linewidth]{example-image-b}
    \caption{Stage 2}\label{fig:s2}
  \end{subfigure}\hfill
  \begin{subfigure}[b]{0.3\textwidth}
    \centering\includegraphics[width=\linewidth]{example-image-c}
    \caption{Stage 3}\label{fig:s3}
  \end{subfigure}

  \bigskip

  \begin{subfigure}[b]{0.3\textwidth}
    \centering\includegraphics[width=\linewidth]{example-image-c}
    \caption{Stage 4}\label{fig:s4}
  \end{subfigure}\hfill
  \begin{subfigure}[b]{0.3\textwidth}
    \centering\includegraphics[width=\linewidth]{example-image-a}
    \caption{Stage 5}\label{fig:s5}
  \end{subfigure}\hfill
  \begin{subfigure}[b]{0.3\textwidth}
    \centering\includegraphics[width=\linewidth]{example-image-b}
    \caption{Stage 6}\label{fig:s6}
  \end{subfigure}
  \caption{Six stages of the assembly process.}
  \label{fig:stages}
\end{figure}

% ----------------------------------------------------------------------
\subsection*{Exercise 8 --- Sizing, framing and cropping}
\exinfo{$\star\star\star$}{8 minutes}{\texttt{scale}, \texttt{height},
\texttt{keepaspectratio}, \texttt{fbox}, \texttt{trim}+\texttt{clip},
\texttt{resizebox}}

\begin{figure}[H]
  \centering
  \begin{subfigure}[b]{0.22\textwidth}
    \centering
    \fbox{\includegraphics[width=\linewidth]{example-image-a}}
    \caption{Framed}\label{fig:framed}
  \end{subfigure}\hfill
  \begin{subfigure}[b]{0.22\textwidth}
    \centering
    \includegraphics[height=2.2cm, keepaspectratio]{example-image-b}
    \caption{Fixed height}\label{fig:height}
  \end{subfigure}\hfill
  \begin{subfigure}[b]{0.22\textwidth}
    \centering
    \includegraphics[width=\linewidth, trim=2cm 1cm 2cm 1cm, clip]{example-image-c}
    \caption{Cropped}\label{fig:crop}
  \end{subfigure}\hfill
  \begin{subfigure}[b]{0.22\textwidth}
    \centering
    \resizebox{\linewidth}{2.2cm}{\includegraphics{example-image-a}}
    \caption{Stretched}\label{fig:stretch}
  \end{subfigure}
  \caption{Four ways to control the size and extent of an image.}
  \label{fig:sizing}
\end{figure}

% ----------------------------------------------------------------------
\newpage
\subsection*{Exercise 9 --- Figure beside a table}
\exinfo{$\star\star\star\star$}{10 minutes}{mixing float types in
\texttt{minipage}s, \texttt{captionof}, independent numbering}

\begin{figure}[H]
\centering
\begin{minipage}[b]{0.46\textwidth}
  \centering
  \includegraphics[width=\linewidth]{example-image-golden}
  \caption{Measured step response.}
  \label{fig:beside}
\end{minipage}\hfill
\begin{minipage}[b]{0.46\textwidth}
  \centering
  \captionof{table}{Extracted parameters.}
  \label{tab:beside}
  \renewcommand{\arraystretch}{1.3}
  \begin{tabular}{|l|c|}
    \hline
    \textbf{Quantity} & \textbf{Value} \\ \hline
    Time constant $\tau$ & $10.0$ ms \\ \hline
    Rise time $t_r$      & $22.0$ ms \\ \hline
    Overshoot            & $4.3\%$   \\ \hline
  \end{tabular}
\end{minipage}
\end{figure}

Figure~\ref{fig:beside} and Table~\ref{tab:beside} describe the same
measurement.

% ----------------------------------------------------------------------
\subsection*{Exercise 10 --- Full exam simulation}
\exinfo{$\star\star\star\star\star$}{15 minutes}{everything at once ---
a wrapped figure, a grouped multi-panel float, sub-references, a footnote
and two cross-references}

\begin{wrapfigure}{l}{0.34\textwidth}
  \centering
  \includegraphics[width=\linewidth]{example-image-16x9}
  \caption{Test rig layout.}
  \label{fig:rig}
\end{wrapfigure}

The apparatus of Figure~\ref{fig:rig} was used to record every trace shown
in Figure~\ref{fig:results}.\footnote{All traces recorded at $25^{\circ}$C
with a $5$\,V step input.} \filler

\begin{figure}[H]
  \centering

  {\large\textbf{Low frequency}}\par\medskip
  \begin{subfigure}[b]{0.28\textwidth}
    \centering\includegraphics[width=\linewidth]{example-image-a}
    \caption{$10$ Hz}\label{fig:f10}
  \end{subfigure}\hfill
  \begin{subfigure}[b]{0.28\textwidth}
    \centering\includegraphics[width=\linewidth]{example-image-b}
    \caption{$50$ Hz}\label{fig:f50}
  \end{subfigure}\hfill
  \begin{subfigure}[b]{0.28\textwidth}
    \centering\includegraphics[width=\linewidth]{example-image-c}
    \caption{$100$ Hz}\label{fig:f100}
  \end{subfigure}

  \bigskip
  {\large\textbf{High frequency}}\par\medskip
  \begin{subfigure}[b]{0.28\textwidth}
    \centering\includegraphics[width=\linewidth]{example-image-c}
    \caption{$1$ kHz}\label{fig:f1k}
  \end{subfigure}\hfill
  \begin{subfigure}[b]{0.28\textwidth}
    \centering\includegraphics[width=\linewidth]{example-image-a}
    \caption{$5$ kHz}\label{fig:f5k}
  \end{subfigure}\hfill
  \begin{subfigure}[b]{0.28\textwidth}
    \centering\includegraphics[width=\linewidth]{example-image-b}
    \caption{$10$ kHz}\label{fig:f10k}
  \end{subfigure}

  \caption{Frequency response of the filter stage, grouped by decade.}
  \label{fig:results}
\end{figure}

Attenuation becomes visible from panel~\subref{fig:f1k} onward, and
Figure~\ref{fig:f10k} shows the roll-off clearly.

\clearpage
% ======================================================================
\part*{PART II --- Solutions and Marking Notes}
\addcontentsline{toc}{section}{PART II --- Solutions and Marking Notes}
\section*{Do not read until you have attempted Part I}
\addcontentsline{toc}{subsection}{Solutions 1--10}

% ----------------------------------------------------------------------
\subsection*{Solution 1}
\begin{lstlisting}[style=tex]
Figure~\ref{fig:warm} shows the reference waveform used throughout.

\begin{figure}[htbp]
  \centering
  \includegraphics[width=0.45\textwidth]{example-image-a}
  \caption{The reference waveform.}
  \label{fig:warm}
\end{figure}

As shown in Figure~\ref{fig:warm}, the response settles quickly.
\end{lstlisting}
\textbf{Marking notes.} Four things must appear in this exact order:
\verb|\centering|, \verb|\includegraphics|, \verb|\caption|, \verb|\label|.
For \emph{figures} the caption goes \textbf{below} the image (for tables it goes
above) --- graders check this. \verb|\label| must follow \verb|\caption|, never
precede it, or \verb|\ref| picks up the wrong counter. Use \verb|~| in
\verb|Figure~\ref{...}| so the number never wraps to the next line alone.
\verb|[htbp]| is the normal placement request; this file uses \verb|[H]| only so
targets appear exactly where printed.

% ----------------------------------------------------------------------
\subsection*{Solution 2}
\begin{lstlisting}[style=tex]
\begin{figure}[htbp]
  \centering
  \begin{minipage}{0.45\textwidth}
    \centering
    \includegraphics[width=\linewidth]{example-image-a}
  \end{minipage}\hfill
  \begin{minipage}{0.45\textwidth}
    \centering
    \includegraphics[width=\linewidth]{example-image-b}
  \end{minipage}
  \caption{Input signal (left) and filtered output (right).}
  \label{fig:pair}
\end{figure}
\end{lstlisting}
\textbf{Marking notes.} Use \texttt{minipage} when the panels need \emph{no}
individual captions --- it is lighter than \texttt{subfigure}. Widths must sum
to less than $1.0$: $0.45 + 0.45 = 0.9$ leaves room for \verb|\hfill| to push
them apart. Inside a minipage, \verb|\linewidth| means \emph{the minipage's}
width, not the page's --- this is the distinction that decides whether your
panels fit. Do not leave a blank line between the two minipages or they will
stack vertically instead of sitting side by side.

% ----------------------------------------------------------------------
\subsection*{Solution 3}
\begin{lstlisting}[style=tex]
\begin{figure}[htbp]
  \centering
  \begin{subfigure}[b]{0.3\textwidth}
    \centering
    \includegraphics[width=\linewidth]{example-image-a}
    \caption{Under-damped}
    \label{fig:under}
  \end{subfigure}\hfill
  \begin{subfigure}[b]{0.3\textwidth}
    \centering
    \includegraphics[width=\linewidth]{example-image-b}
    \caption{Critically damped}
    \label{fig:crit}
  \end{subfigure}\hfill
  \begin{subfigure}[b]{0.3\textwidth}
    \centering
    \includegraphics[width=\linewidth]{example-image-c}
    \caption{Over-damped}
    \label{fig:over}
  \end{subfigure}
  \caption{Three damping regimes of a second-order system.}
  \label{fig:damping}
\end{figure}

Panel~\subref{fig:crit} of Figure~\ref{fig:damping} settles fastest, while
Figure~\ref{fig:under} overshoots.
\end{lstlisting}
\textbf{Marking notes.} Each \texttt{subfigure} has its own
\verb|\caption|+\verb|\label|; the float has one more of each at the end.
Note the difference in output: \verb|\ref{fig:under}| prints the \emph{full}
number \ref{fig:under} (figure number plus sub-letter) while
\verb|\subref{fig:crit}| prints only \subref{fig:crit}. The optional
\verb|[b]| bottom-aligns panels so their captions line up even when images
differ in height. Three panels at $0.3$ leaves $0.1$ for the two \verb|\hfill|s.

% ----------------------------------------------------------------------
\subsection*{Solution 4}
\begin{lstlisting}[style=tex]
\begin{figure}[htbp]
  \centering

  {\large\textbf{Star}}\par\medskip
  \begin{subfigure}[b]{0.3\textwidth}
    \centering
    \includegraphics[width=0.6\linewidth]{example-image-a}
    \caption{The Sun}\label{fig:sun}
  \end{subfigure}

  \bigskip
  {\large\textbf{Inner Planets}}\par\medskip
  \begin{subfigure}[b]{0.22\textwidth}
    \centering\includegraphics[width=\linewidth]{example-image-b}
    \caption{Mercury}\label{fig:mercury}
  \end{subfigure}\hfill
  ... three more at 0.22\textwidth ...

  \bigskip
  {\large\textbf{Outer Planets}}\par\medskip
  ... four more at 0.22\textwidth ...

  \caption{The star and the eight planets, arranged into the inner and
  outer groups. The displayed sizes and distances are not to scale.}
  \label{fig:grid}
\end{figure}
\end{lstlisting}
\textbf{Marking notes.} This is the most time-expensive item in either paper, so
the technique matters more than the typing. \textbf{Write one subfigure block,
compile it, then copy--paste it eight times} and change only three things per
copy: the image name, the caption word, the label. Never retype from scratch.

The group headings are ordinary centred bold text \emph{inside} the float:
\verb|{\large\textbf{Star}}\par\medskip|. The \verb|\par| is essential --- it
closes the paragraph so the heading sits on its own line; without it the
heading and the first panel collide. A \textbf{blank line} between panel rows
is what breaks to a new row; \verb|\bigskip| just adds the vertical gap. Four
panels at $0.22$ sum to $0.88$, leaving room for three \verb|\hfill|s.

% ----------------------------------------------------------------------
\subsection*{Solution 5}
\begin{lstlisting}[style=tex]
\begin{wrapfigure}{r}{0.38\textwidth}   % r = right, l = left
  \centering
  \includegraphics[width=\linewidth]{example-image-a}
  \caption{A wrapped figure on the right.}
  \label{fig:wrap}
\end{wrapfigure}

Figure~\ref{fig:wrap} sits on the right while this paragraph flows
around it. Lorem ipsum ... (enough text to fill the figure's height)
\end{lstlisting}
\textbf{Marking notes.} \texttt{wrapfigure} takes \emph{two} mandatory
arguments: the side (\texttt{r}, \texttt{l}, or capital \texttt{R}/\texttt{L} to
allow floating) and the reserved width. It is \textbf{not} a float in the usual
sense --- place it immediately before the paragraph that should wrap, and never
inside a list, another float, or directly after \verb|\section|, where it
misbehaves. You need enough following text to fill the figure's height, or the
wrap spills untidily onto the next paragraph. If the wrap looks broken, add
\verb|\bigskip| before it or shift it one paragraph earlier.

% ----------------------------------------------------------------------
\subsection*{Solution 6}
\begin{lstlisting}[style=tex]
% inside the middle subfigure:
\includegraphics[width=0.7\linewidth, angle=90]{example-image-a}
% inside the third subfigure:
\rotatebox{30}{\fbox{Rotated text}}
% whole-float rotation (package rotating), used on its own:
\begin{sidewaysfigure}
  \centering
  \includegraphics[width=0.8\textwidth]{example-image-a}
  \caption{A landscape figure on its own page.}
  \label{fig:sideways}
\end{sidewaysfigure}
\end{lstlisting}
\textbf{Marking notes.} Three distinct tools. \verb|angle=| inside
\verb|\includegraphics| rotates \emph{the image}, and it interacts with
\verb|width=|: the width is applied \emph{before} rotation, so a rotated image
ends up as wide as the original was tall --- always check the result rather than
assuming. \verb|\rotatebox{deg}{...}| (from \texttt{graphicx}) rotates
\emph{anything}, including text and tables. \texttt{sidewaysfigure} (from
\texttt{rotating}) turns the entire float, caption included, onto its own
landscape page --- reach for it when a wide figure will not fit upright.
Positive angles rotate counter-clockwise.

% ----------------------------------------------------------------------
\subsection*{Solution 7}
\begin{lstlisting}[style=tex]
\begin{figure}[htbp]
  \centering
  \begin{subfigure}[b]{0.3\textwidth}
    \centering\includegraphics[width=\linewidth]{example-image-a}
    \caption{Stage 1}\label{fig:s1}
  \end{subfigure}\hfill
  \begin{subfigure}[b]{0.3\textwidth} ... \end{subfigure}\hfill
  \begin{subfigure}[b]{0.3\textwidth} ... \end{subfigure}

  \bigskip                       % <-- blank line above = new row

  \begin{subfigure}[b]{0.3\textwidth} ... \end{subfigure}\hfill
  \begin{subfigure}[b]{0.3\textwidth} ... \end{subfigure}\hfill
  \begin{subfigure}[b]{0.3\textwidth} ... \end{subfigure}
  \caption{Six stages of the assembly process.}
  \label{fig:stages}
\end{figure}
\end{lstlisting}
\textbf{Marking notes.} The row break is the \emph{blank line}, not
\verb|\bigskip| --- \verb|\bigskip| only adds space. Equivalently you can end a
row with \verb|\\|. Keep every panel the same declared width or the columns will
not align between rows. Sub-labels run (a)--(f) continuously across both rows;
they are numbered in source order, so if you reorder panels the letters follow.

% ----------------------------------------------------------------------
\subsection*{Solution 8}
\begin{lstlisting}[style=tex]
\fbox{\includegraphics[width=\linewidth]{example-image-a}}          % framed
\includegraphics[height=2.2cm, keepaspectratio]{example-image-b}    % fixed height
\includegraphics[width=\linewidth, trim=2cm 1cm 2cm 1cm, clip]{example-image-c}
\resizebox{\linewidth}{2.2cm}{\includegraphics{example-image-a}}    % stretched
% also valid:
\includegraphics[scale=0.4]{example-image}       % fraction of natural size
\end{lstlisting}
\textbf{Marking notes.} \verb|scale=| is relative to the image's natural size;
\verb|width=|/\verb|height=| are absolute. Giving \emph{both} width and height
distorts the image unless you add \verb|keepaspectratio|. \verb|trim=| takes
four lengths in the order \textbf{left bottom right top}, and it does nothing
visible unless you also pass \verb|clip| --- forgetting \verb|clip| is the
classic bug. \verb|\resizebox{W}{H}{...}| forces content into an exact box
(use \verb|!| for one dimension to preserve the ratio, e.g.
\verb|\resizebox{5cm}{!}{...}|); it works on tables too, which is the quickest
rescue for a table that overflows the margin.

% ----------------------------------------------------------------------
\subsection*{Solution 9}
\begin{lstlisting}[style=tex]
\begin{figure}[htbp]
\centering
\begin{minipage}[b]{0.46\textwidth}
  \centering
  \includegraphics[width=\linewidth]{example-image-golden}
  \caption{Measured step response.}
  \label{fig:beside}
\end{minipage}\hfill
\begin{minipage}[b]{0.46\textwidth}
  \centering
  \captionof{table}{Extracted parameters.}
  \label{tab:beside}
  \begin{tabular}{|l|c|}
    \hline \textbf{Quantity} & \textbf{Value} \\ \hline
    Time constant $\tau$ & $10.0$ ms \\ \hline
    Rise time $t_r$      & $22.0$ ms \\ \hline
    Overshoot            & $4.3\%$   \\ \hline
  \end{tabular}
\end{minipage}
\end{figure}
\end{lstlisting}
\textbf{Marking notes.} A \texttt{table} float cannot be nested inside a
\texttt{figure} float, so the trick is to put \emph{both} in minipages and use
\verb|\captionof{table}{...}| (from \texttt{caption}, which \texttt{subcaption}
already loads) to borrow the table counter. The figure keeps the figure counter,
the table keeps the table counter, and both \verb|\ref|s work independently.
Remember the table caption goes \textbf{above} its tabular while the figure
caption goes \textbf{below} its image --- that asymmetry is visible here in one
float and is a favourite examiner detail. Escape the percent sign as
\verb|4.3\%|.

% ----------------------------------------------------------------------
\subsection*{Solution 10}
\begin{lstlisting}[style=tex]
\begin{wrapfigure}{l}{0.34\textwidth}
  \centering
  \includegraphics[width=\linewidth]{example-image-16x9}
  \caption{Test rig layout.}\label{fig:rig}
\end{wrapfigure}

The apparatus of Figure~\ref{fig:rig} was used to record every trace shown
in Figure~\ref{fig:results}.\footnote{All traces recorded at $25^{\circ}$C
with a $5$\,V step input.} ... enough text to fill the wrap ...

\begin{figure}[htbp]
  \centering
  {\large\textbf{Low frequency}}\par\medskip
  \begin{subfigure}[b]{0.28\textwidth}
    \centering\includegraphics[width=\linewidth]{example-image-a}
    \caption{$10$ Hz}\label{fig:f10}
  \end{subfigure}\hfill
  ... two more ...

  \bigskip
  {\large\textbf{High frequency}}\par\medskip
  ... three more ...

  \caption{Frequency response of the filter stage, grouped by decade.}
  \label{fig:results}
\end{figure}

Attenuation becomes visible from panel~\subref{fig:f1k} onward, and
Figure~\ref{fig:f10k} shows the roll-off clearly.
\end{lstlisting}
\textbf{Marking notes.} Six things are graded at once. The \verb|\footnote|
here works because it sits in \emph{body text}, not inside the float --- had it
been inside, you would need \verb|\footnotemark| plus \verb|\footnotetext|.
Units and degrees inside captions are ordinary maths: \verb|$10$ Hz|,
\verb|$25^{\circ}$C|, and \verb|$5$\,V| uses a thin space before the unit.
Build order under time pressure: place the wrapped figure first, then one
subfigure block, compile, then copy it five times. Finish with the two
cross-references, and compile \textbf{twice} so every \verb|\ref| resolves.

% ======================================================================
\clearpage
\section*{Drill Routine}
\addcontentsline{toc}{section}{Drill Routine}

\begin{enumerate}[leftmargin=*]
\item \textbf{Round 1 --- untimed.} Work through 1, 2, 3, 5 with Part II open.
  The aim is recognition of the four base patterns: single float, minipage pair,
  subfigure row, wrapped float.
\item \textbf{Round 2 --- timed, base patterns.} Close Part II. Reproduce
  1, 2, 3, 6 in under 20 minutes total.
\item \textbf{Round 3 --- the grid.} Exercise 4 alone, in 12 minutes, using
  copy--paste discipline. This is the highest-value drill in the pack: the same
  structure appeared as a ten-mark item in the real paper.
\item \textbf{Round 4 --- awkward layouts.} Exercises 7, 8, 9 in 25 minutes.
\item \textbf{Round 5 --- exam pace.} Exercise 10 cold in 15 minutes, including
  the surrounding prose, the footnote and both cross-references.
\item \textbf{Round 6 --- from memory.} Write Exercise 4 with nothing open. When
  the nine-panel grid costs you under ten minutes, figures are solved.
\end{enumerate}

\section*{The Figure Checklist}
\addcontentsline{toc}{section}{The Figure Checklist}

Run this over every float before moving on:
\begin{enumerate}[label=\textbf{\arabic*.}, leftmargin=*]
\item \verb|\centering| is present (not the \texttt{center} environment, which
  adds unwanted vertical space inside a float).
\item The caption is \textbf{below} the image, and \verb|\label| comes
  \textbf{after} \verb|\caption|.
\item Every panel width is declared and the widths sum to less than $1.0$.
\item Panels are separated by \verb|\hfill| with \textbf{no blank line} between
  them; a blank line is used only to start a new row.
\item Each label is unique and every one is actually referenced in the text.
\item The document has been compiled \textbf{twice}.
\end{enumerate}

\section*{The Six Errors That Cost the Most Marks}
\addcontentsline{toc}{section}{The Six Errors That Cost the Most Marks}

\renewcommand{\arraystretch}{1.4}
\begin{tabularx}{\textwidth}{|c|l|X|}
\hline
\textbf{\#} & \textbf{Mistake} & \textbf{Symptom and fix} \\
\hline
1 & Blank line between panels &
  Panels stack vertically instead of sitting side by side. Remove the empty
  line; keep only \verb|\hfill|. \\
\hline
2 & \verb|\label| before \verb|\caption| &
  \verb|\ref| prints the number of the \emph{previous} float. Always
  \verb|\caption| then \verb|\label|. \\
\hline
3 & Panel widths sum to $\geq 1.0$ &
  \texttt{Overfull \textbackslash hbox} and panels spill into the margin.
  Three panels: use $0.3$; four: use $0.22$. \\
\hline
4 & \verb|\textwidth| inside a minipage &
  The image bursts out of its panel. Inside a minipage or subfigure use
  \verb|\linewidth|. \\
\hline
5 & \verb|trim| without \verb|clip| &
  Nothing appears to crop. Both keys are required. \\
\hline
6 & Compiling once &
  \verb|\ref| shows \texttt{??} and sub-references are blank. Run
  \texttt{pdflatex} twice. \\
\hline
\end{tabularx}

\vfill
\begin{center}\small\itshape
One panel, compiled and correct, then copy it eight times.
\end{center}

\end{document}
